\RequirePackage[orthodox]{nag}
\documentclass[11pt]{article}

%% Define the include path
\makeatletter
\providecommand*{\input@path}{}
\g@addto@macro\input@path{{include/}{../../include/}}
\makeatother

\usepackage{../../include/akazachk}

\title{Software Engineering for Systems Modeling}
\author{Andres Espinosa}
\begin{document}
\pgfplotsset{compat=1.18}
\maketitle

\tableofcontents
\section{Tuesday 03/24/2026}

\subsection{V-Model}

Requirements Analysis -> High level design -> Detailed specifications -> Coding -> Unit testing -> Integration testing -> Operational testing

Requirements can be divided into that of problem (stakeholder) and solution (system) requirements.
The requirements of the problem and stakeholder are concerned with the satisfaction of the output of the system.
It is concerned with what a stakeholder would like to have from the system and how they might interact with it.
The requirements of the solution and system are those that lie in the bounds of the system and define what it does.

\subsection{MBSE}
\textbf{Model-based Systems Engineering} defines a standard comprehensive perspect of a system and develops a single model for the process.
It lowers the cost of changes throughout the process and establishes clear boundaries of both components of the system and the processes associated.

\subsection{SysML}
SysML shares overlap with UML, a modeling language and standard that provides a visual standard for components of software that interact with each other.
With a modeling language, we can construct concrete and visual representations of previously abstract ideas.
A system could be represented as a list of documents describing everything in detail.
However, a model enables visual representations more flexible and comprehensable.

SysML supports the specification, analysis, design, verification, and validation of systems that include hardware, software, data, personel, procedures, and facilities.
It is a visual modeling language that provides semantics and notations for models, which at times, can be quite tedious and has inhibited its prevalence.

The four pillars of SysML include: 
\begin{enumerate}
    \item Structure
    \item Behavior
    \item Requirements
    \item Parametrics
\end{enumerate}

A SysML drawing has a few components
\begin{itemize}
    \item Analysis diagram
    \item Internal block diagram
    \item Package diagram
    \item Parametric diagram
    \item Requirement diagram
    \item Sequence diagram
    \item State machine diagram
    \item Use case diagram
    \item Allocation tables
\end{itemize}

Some important questions for defining a system:
\begin{itemize}
    \item What is its purpose and how is its purpose accomplished?
    \item How does it interact with its environment?
    \item Who is the user and how do they use it?
    \item What are the components necessary to make it function, how do they interoperate and how do they interact with the environment?
    \item What are the capabilities of the system and how are they tested?
    \item What is the cost of the system?
    \item What is the lifespan of the system?
\end{itemize}

\end{document}